% Standalone proof. Shared notation and prose conventions: analyses/README.md.
\documentclass[11pt,a4paper]{article}
\usepackage[T1]{fontenc}
\usepackage[utf8]{inputenc}
\usepackage{amsmath,amsthm}
\usepackage{newtxtext,newtxmath}
\usepackage[a4paper,margin=21mm,headheight=15pt]{geometry}
\usepackage{microtype,booktabs,array,tabularx,longtable}
\usepackage[table]{xcolor}
\usepackage{tikz,fancyhdr,enumitem,seqsplit}
\definecolor{Teal}{HTML}{176B70}
\definecolor{Ink}{HTML}{203443}
\definecolor{Pale}{HTML}{EDF5F5}
\definecolor{Line}{HTML}{C7D8DB}
\definecolor{Muted}{HTML}{536974}
\definecolor{Warm}{HTML}{F8F4EC}
\usepackage[colorlinks=true,linkcolor=Teal,urlcolor=Teal]{hyperref}
\color{Ink}
\newcommand{\blank}{\lambda}
\newcommand{\digits}{\{0,1,\ldots,9\}}
\newcommand{\ph}[1]{\mathsf{#1}}
\newcommand{\EXPERT}{\mathcal E}
\newcommand{\CTRL}{\mathcal C}
\newcommand{\PTO}{\operatorname{PTO}}
\newcommand{\triple}[3]{(#1,#2,\mathsf{#3})}
\newcommand{\code}[1]{\texttt{#1}}
\newcommand{\hash}[1]{{\scriptsize\ttfamily\seqsplit{#1}}}
\newtheorem{theorem}{Theorem}
\newtheorem{lemma}{Lemma}
\newtheorem{corollary}{Corollary}
\setlength{\parindent}{0pt}
\setlength{\parskip}{6pt}
\setlist[itemize]{nosep,leftmargin=1.2em}
\pagestyle{fancy}
\fancyhf{}
\fancyhead[L]{\small\sffamily\color{Muted}Decimal Multiplication}
\fancyhead[R]{\small\sffamily\color{Muted}Metric Extraction}
\fancyfoot[L]{\scriptsize\sffamily\color{Muted}Exact PTO Equivalence}
\fancyfoot[R]{\small\thepage}
\renewcommand{\headrulewidth}{.35pt}
\renewcommand{\footrulewidth}{0pt}
\newcommand{\sectiontitle}[2]{\refstepcounter{section}\vspace*{2pt}{\LARGE\bfseries\color{Teal}\thesection\quad #2}\par\vspace{7pt}}

\date{}
\pdfinfoomitdate=1
\pdftrailerid{}
\begin{document}
{\Huge\bfseries Exact PTO Equivalence\\[3pt]Decimal Multiplication}\par
\vspace{10pt}
{\Large Metric-Extracted Controller}\par
{\small\color{Muted}\path{analyses/metric_extraction/DFST/mult10/warmup_0_steps_4000000/}}\par
\vspace{7pt}

\vspace{4pt}
\sectiontitle{1 \enspace STATEMENT AND SCOPE}{Statement and Scope}\label{sec:common1}
\textbf{Controllers.}
The export \path{minimized_partial.json} defines controller $\CTRL$, with
$182$ states, initial state $0$, and halt state $181$.
The function \code{mult\_grid(10)} defines expert $\EXPERT$, with
$108$ states, initial state $0$, and halt state $107$.
The adjacent expert export matches this function entry for entry.
All state numbers below follow the original exports.

\textbf{Input Domain.}
Let $D=\digits$. Row $0$ initially contains the contiguous word $a\,x\,b$,
starting at column $0$, where $a\in D^m$, $b\in D^n$, and $m,n\ge1$.
All remaining grid cells contain blanks. The head starts at $(0,0)$.
Write $A=\operatorname{val}_{10}(a)$ and $B=\operatorname{val}_{10}(b)$.
The domain permits arbitrary finite operand lengths, unequal lengths,
leading zeros, and zero-valued operands.

\textbf{Execution Convention.}
The grid uses $\mathbb Z^2$, with coordinates increasing rightward and
downward. The symbol $\blank$ denotes a blank. A transition $(q',w,M)$
writes $w$, performs motion $M$, and enters state $q'$.
Motions $L,R,D,U$ move one cell; motion $S$ writes and stops without moving.
An undefined transition would prevent continuation.

\textbf{Policy-Trajectory Observations.}
Define the PTO on input $a\,x\,b$ by
\[
\PTO(a\,x\,b)=\bigl((r_t,w_t,M_t)\bigr)_{t=0}^{T-1},
\]
where $r_t,w_t,M_t$ denote the read symbol, written symbol, and motion
at step $t$. Count the final stopping action exactly once, without halt-loop padding.

\begin{theorem}[Exact PTO Equivalence]\label{thm:main}
For every input in the stated domain,
\[
\PTO_{\CTRL}(a\,x\,b)=\PTO_{\EXPERT}(a\,x\,b).
\]
Both executions terminate on the same step. The complete grids and head
positions agree at every step. No undefined controller entry occurs.
The final grid contains only the $(m+n)$-digit decimal representation of
$AB$, with leading-zero padding, on row $n-1$ in columns $-n,\ldots,m-1$.
The final head stops at $(-n-1,n-1)$.
\end{theorem}

The theorem concerns the exported discrete controller. The proof combines
geometric and arithmetic invariants with a finite simulation certificate;
no bound on $m$ or $n$ enters the argument.
\clearpage
\sectiontitle{2 \enspace EXPERT BEHAVIOR}{Expert Behavior}\label{sec:common2}
\subsection{Copying and Partial Products}
Write $b=b_{n-1}\cdots b_0$, so $B=\sum_{i=0}^{n-1}b_i10^i$.
Assign global weight $10^{m-1-z}$ to column $z$.
Column $m-1$ carries the units digit of the final product.

\begin{lemma}[Copy Staircase]\label{lem:copy}
Before product formation, the expert constructs exactly $n$ occupied rows.
For $0\le i<n$, row $i$ contains
\[
a\,x\,b_{n-1}\cdots b_i
\quad\text{starting in column }-i. \tag{1}
\]
All remaining cells contain blanks.
\end{lemma}
\begin{proof}
Row $0$ contains the input. On row $i$, the expert scans to the right blank,
steps left over $b_i$, and inspects the preceding cell. A digit at this
position implies at least two remaining factor digits. The expert copies
every symbol except $b_i$ one row down and one column left, from right to left.
Each destination contains a blank before writing; the source row remains unchanged.

The first descent reaches a blank. Subsequent descents reach the copied
right neighbor: a digit, except when copying the last digit of $a$, where
the lower cell contains the copied $x$. The descent from the source $x$
reaches a digit. After copying the first digit of $a$, the expert reads the
left blank and descends onto the first digit of the new row.
The new row satisfies (1), with one fewer factor digit.

When one factor digit remains, the second left step reaches $x$ and starts
product formation. Thus $n-1$ copies produce exactly the stated staircase.
\end{proof}

\begin{lemma}[Aligned Partial Products]\label{lem:products}
After product formation, row $i$ contains the $(m+1)$-digit representation
of $A b_i$ on columns $-i-1,\ldots,m-i-1$.
The leading digit belongs to $\{0,\ldots,8\}$.
In global column weights, this row represents $A b_i10^i$.
All remaining cells contain blanks.
\end{lemma}
\begin{proof}
The expert processes rows from $n-1$ upwards. On row $i$, the expert erases
the selected factor $f=b_i$, the remaining factor prefix, and $x$, then
walks left across $a$. For $f=0$, each digit becomes $0$, followed by one
padding $0$ on the left blank. For $1\le f\le9$, the control state stores
carry $c$, initially $0$, and applies
\[
fd+c=w+10c',\qquad
w=(fd+c)\bmod10,\quad c'=\lfloor(fd+c)/10\rfloor. \tag{2}
\]
The bound $0\le c<f$ persists because $fd+c\le9f+f-1=10f-1$.
The usual suffix invariant for (2) gives the product digits; the left blank
receives the final carry. Thus the leading digit never exceeds $f-1\le8$.

The right scan ends at column $m-i$ and moves up. For $i>0$, the upper cell
contains the last digit of $a$, followed by $x$ and the remaining factor prefix.
The expert scans this prefix, selects the last digit, and repeats.
Every higher row contains at least two factor digits, so at least one further
digit follows the selected factor during leftward erasure.
For $i=0$, the upper cell contains a blank, starting summation on row $0$.
\end{proof}

Every scan crosses a finite word. Copying shortens the factor prefix;
product formation visits each occupied row once. Both passes terminate.

\newpage
\subsection{Accumulation and Carry Confinement}
\begin{lemma}[Fixed-Width Accumulator]\label{lem:sum}
After merging through row $k$, $0\le k<n$, all rows above $k$ contain blanks.
Row $k$ contains a contiguous $(m+k+1)$-digit accumulator on columns
$-k-1,\ldots,m-1$, with value
\[
S_k=A\sum_{i=0}^{k}b_i10^i.
\]
Rows below $k$ retain the partial products from Lemma~\ref{lem:products}.
At inspection of the next row, expert state $88$ has head $(-k-2,k+1)$.
\end{lemma}
\begin{proof}
For $k=0$, the expert scans left across the product on row $0$, reaches
column $-2$, and descends to inspect row $1$.
If row $k+1$ exists, inspection reads the fresh leading digit in
$\{0,\ldots,8\}$. The expert scans the lower product to the right blank at
$m-k-1$ and moves up. The upper cell contains an accumulator digit, since
\[
-k-1\le m-k-1\le m-1.
\]
The expert scans to the upper right blank and steps left to column $m-1$.
The lower product ends at $m-k-2$, so the first lower cell contains a blank.

For upper digit $u$, lower digit $v$ (a blank contributes zero), and incoming
carry $c\in\{0,1\}$, the update applies
\[
u+v+c=w+10c',\qquad
w\in D,\quad c'\in\{0,1\}. \tag{3}
\]
The expert erases $u$, writes $w$ below, returns to the erased upper cell,
and moves left. The processed lower suffix contains the correct sum suffix;
the unprocessed cells retain the old digits; $c'$ records the contribution
to the next column. Equation~(3) preserves this invariant.

The lower word extends one column farther left than the upper word.
After the upper word ends, carry zero preserves this lower leading digit.
Carry one increments the leading digit and clears, since the digit never
exceeds $8$. The expert reaches column $-k-3$ on row $k+1$ and descends.
Row $k$ contains only blanks; row $k+1$ represents
$S_k+A b_{k+1}10^{k+1}=S_{k+1}$ on the claimed support.
\end{proof}

\begin{lemma}[Boundary Restrictions]\label{lem:boundary}
A carry-one update beneath an upper digit reads an original lower digit.
A blank lower update guarantees another upper digit on the next left step.
A carry beyond the upper word reads a lower digit in $\{0,\ldots,8\}$ and clears.
\end{lemma}
\begin{proof}
The blank gap lies strictly to the right of the original lower word.
The first update has carry zero, and $u+0\le9$, so no carry originates
inside this gap. Carry one can first arise only inside the lower digit block;
further leftward updates stay inside that block until the leading digit.

A blank lower cell has column $z\ge m-k-1$. The next upper column satisfies
$z-1\ge m-k-2\ge-k-1$, using $m\ge1$, and therefore still contains a digit.
Finally, the lower leading digit lies one column left of the upper word.
Lemma~\ref{lem:products} bounds this untouched digit by $8$; adding $1$ clears
any final carry without extending the fixed-width support.
\end{proof}

Each merge consumes finitely many upper digits and at most one extra carry
column. After $n-1$ merges, state $88$ reads the blank below the final row,
moves up to state $106$, and emits $S$ into halt state $107$.
Lemma~\ref{lem:sum} gives the final product, padding, and head position.

\clearpage
\sectiontitle{3 \enspace OBSERVATION PROTOCOL}{Observation Protocol}\label{sec:common3}
Section~\ref{sec:common4} gives $207$ phases. Phase $i$ pairs expert state $g_i$ with
controller state $q\in Q_i$, on an identical grid and head.
An edge $i\xrightarrow{x}j$ permits read symbol $x$ and advances to phase $j$.
The protocol begins at phase $0$ and ends at phase $206$.
All phase sets contain one controller state except
$Q_{201}=\{26,106\}$.

\begin{lemma}[Protocol Coverage]\label{lem:coverage}
Every expert execution in the stated domain follows the protocol in Section~\ref{sec:common4}
from phase $0$ to phase $206$.
\end{lemma}
\begin{proof}
The following geometric and arithmetic conditions justify every read restriction.

\textbf{Copying.}
The initial scan and every restart encounter a nonempty left operand,
then $x$, then a nonempty right operand before the right blank.
The second left step reads a factor digit or $x$.
Lemma~\ref{lem:copy} distinguishes the first descent (blank), ordinary descents
(digits), and the descent directly above the copied operator ($x$).
Every destination contains a blank. After the first copied symbol, the next
source symbol contains a digit or $x$; no left blank can occur yet.
After copying $x$, the next source symbol contains the last digit of $a$.
The ordinary copy loop eventually reaches the left blank and restarts.

\textbf{Products.}
The bottom row contains one factor digit; erasing that digit immediately
exposes $x$. Every higher row contains at least two factor digits;
erasing the selected factor therefore exposes another digit before $x$.
After erasing $x$, the first product update reads a digit because $m\ge1$.
Later updates read digits followed by the left blank.
Equation~(2) determines the next carry phase and preserves the carry bound.
The completed product scan reads digits followed by a blank.
Lemma~\ref{lem:products} places the head on an upper digit followed by $x$
and a nonempty factor prefix, or on a blank above row $0$.

\textbf{Alignment and Summation.}
Summation starts on the right blank of row $0$. The initial left scan crosses
a nonempty product. Inspection below reads a fresh leading digit from
$\{0,\ldots,8\}$ or a blank below the last row.
After a leading digit, at least one lower digit remains because $m\ge1$.
Lemma~\ref{lem:sum} places the first upper alignment step on a digit;
further alignment steps read digits until the right blank.
The first upper units digit has carry zero and a blank lower cell.

After every lower update, the upper return cell contains the erased blank.
A blank lower update leads to a nonblank upper read by
Lemma~\ref{lem:boundary}. An update inside the lower digit block leads to
an upper digit or the left blank. Carry-one digit updates remain inside
the lower digit block. At the end of the upper word, carry zero preserves
the untouched lower leading digit, while carry one increments that digit
and clears. Both paths reach the blank immediately left of the merged word.
The final stop reads a blank.

These conditions give exactly the permitted branches used by the expert.
Induction over individual moves, copied rows, and merged rows establishes
coverage. The finite-pass arguments in Lemmas~\ref{lem:copy}--\ref{lem:sum}
give termination at phase $206$.
\end{proof}

The protocol may permit additional observation sequences. Checking every
permitted branch therefore suffices for all valid grid executions.

\clearpage
\sectiontitle{4 \enspace SIMULATION CERTIFICATE}{Simulation Certificate}\label{sec:common4}
\textbf{Finite Local Certificate.}
For every protocol edge $i\xrightarrow{x}j$ and every $q\in Q_i$,
direct comparison with the two exports gives
\begin{equation}
\delta_\EXPERT(g_i,x)=(g_j,w,M),\qquad
\delta_\CTRL(q,x)=(q',w,M),\qquad q'\in Q_j. \tag{4}
\end{equation}
Every required entry exists. A transition reaches halt exactly when
$j=206$, with motion $S$. The companion checker verifies all $1{,}509$
obligations across $208$ phase--state pairs. These obligations cover all
$1{,}259$ defined nonterminal controller entries. The remaining $12$ defined
entries contain halt loops, which never contribute to a PTO.

\subsection{Phases and Controller-State Relation}
The table gives the phase index $i$, expert state $g_i$, controller set $Q_i$,
and permitted reads with successor phase indices.
For a digit vector $[j_0,\ldots,j_9]$, read $d$ advances to $j_d$.
The notation $D:j$ sends every decimal digit to phase $j$;
$0{:}8:j$ permits only digits $0,\ldots,8$.
Unlisted reads never occur in the corresponding phase of a valid input.

The expert transition supplies the write and motion for each edge;
Equation~(4) checks both exports against that action. The copying moves
follow Lemma~\ref{lem:copy}; product updates follow (2); addition updates
follow (3). Phase labels describe positions before the next action.
\begingroup
\fontsize{8.7}{10.5}\selectfont
\setlength{\tabcolsep}{3pt}
\renewcommand{\arraystretch}{1.17}
\begin{longtable}{r p{4.0cm} r l p{8.1cm}}
\toprule\rowcolor{Pale}
$i$ & Phase & $g_i$ & $Q_i$ & Read $:$ next phase\\\midrule
\endfirsthead
\multicolumn{5}{l}{\small\sffamily\color{Muted}Simulation Certificate (Continued)}\\[5pt]
\toprule\rowcolor{Pale}
$i$ & Phase & $g_i$ & $Q_i$ & Read $:$ next phase\\\midrule
\endhead
\midrule\multicolumn{5}{r}{\scriptsize Continued on the next page}\\
\endfoot
\bottomrule
\endlastfoot
0 & First left-operand digit & 0 & $\{0\}$ & $D:1$ \\
1 & Scan left operand & 0 & $\{0\}$ & $D:1\quad x:2$ \\
2 & First factor digit & 0 & $\{1\}$ & $D:3$ \\
3 & Scan right operand & 0 & $\{1\}$ & $D:3\quad \blank:4$ \\
4 & Skip last factor digit & 1 & $\{2\}$ & $D:5$ \\
5 & Copy or start products & 2 & $\{3\}$ & $D:[6,8,10,12,14,16,18,20,22,24]$\newline$x:62$ \\
6 & First descent, digit 0 & 3 & $\{172\}$ & $\blank:7$ \\
7 & First copied write, digit 0 & 14 & $\{173\}$ & $\blank:26$ \\
8 & First descent, digit 1 & 4 & $\{123\}$ & $\blank:9$ \\
9 & First copied write, digit 1 & 15 & $\{124\}$ & $\blank:26$ \\
10 & First descent, digit 2 & 5 & $\{64\}$ & $\blank:11$ \\
11 & First copied write, digit 2 & 16 & $\{65\}$ & $\blank:26$ \\
12 & First descent, digit 3 & 6 & $\{4\}$ & $\blank:13$ \\
13 & First copied write, digit 3 & 17 & $\{5\}$ & $\blank:26$ \\
14 & First descent, digit 4 & 7 & $\{87\}$ & $\blank:15$ \\
15 & First copied write, digit 4 & 18 & $\{88\}$ & $\blank:26$ \\
16 & First descent, digit 5 & 8 & $\{39\}$ & $\blank:17$ \\
17 & First copied write, digit 5 & 19 & $\{40\}$ & $\blank:26$ \\
18 & First descent, digit 6 & 9 & $\{163\}$ & $\blank:19$ \\
19 & First copied write, digit 6 & 20 & $\{164\}$ & $\blank:26$ \\
20 & First descent, digit 7 & 10 & $\{51\}$ & $\blank:21$ \\
21 & First copied write, digit 7 & 21 & $\{52\}$ & $\blank:26$ \\
22 & First descent, digit 8 & 11 & $\{128\}$ & $\blank:23$ \\
23 & First copied write, digit 8 & 22 & $\{129\}$ & $\blank:26$ \\
24 & First descent, digit 9 & 12 & $\{79\}$ & $\blank:25$ \\
25 & First copied write, digit 9 & 23 & $\{80\}$ & $\blank:26$ \\
26 & Source after first copy & 25 & $\{6\}$ & $D:[27,29,31,33,35,37,39,41,43,45]$\newline$x:48$ \\
27 & Ordinary descent, digit 0 & 3 & $\{178\}$ & $D:28$ \\
28 & Ordinary copied write, digit 0 & 14 & $\{159\}$ & $\blank:47$ \\
29 & Ordinary descent, digit 1 & 4 & $\{134\}$ & $D:30$ \\
30 & Ordinary copied write, digit 1 & 15 & $\{135\}$ & $\blank:47$ \\
31 & Ordinary descent, digit 2 & 5 & $\{68\}$ & $D:32$ \\
32 & Ordinary copied write, digit 2 & 16 & $\{67\}$ & $\blank:47$ \\
33 & Ordinary descent, digit 3 & 6 & $\{125\}$ & $D:34$ \\
34 & Ordinary copied write, digit 3 & 17 & $\{11\}$ & $\blank:47$ \\
35 & Ordinary descent, digit 4 & 7 & $\{166\}$ & $D:36$ \\
36 & Ordinary copied write, digit 4 & 18 & $\{90\}$ & $\blank:47$ \\
37 & Ordinary descent, digit 5 & 8 & $\{165\}$ & $D:38$ \\
38 & Ordinary copied write, digit 5 & 19 & $\{131\}$ & $\blank:47$ \\
39 & Ordinary descent, digit 6 & 9 & $\{55\}$ & $D:40$ \\
40 & Ordinary copied write, digit 6 & 20 & $\{42\}$ & $\blank:47$ \\
41 & Ordinary descent, digit 7 & 10 & $\{170\}$ & $D:42$ \\
42 & Ordinary copied write, digit 7 & 21 & $\{148\}$ & $\blank:47$ \\
43 & Ordinary descent, digit 8 & 11 & $\{150\}$ & $D:44$ \\
44 & Ordinary copied write, digit 8 & 22 & $\{116\}$ & $\blank:47$ \\
45 & Ordinary descent, digit 9 & 12 & $\{167\}$ & $D:46$ \\
46 & Ordinary copied write, digit 9 & 23 & $\{54\}$ & $\blank:47$ \\
47 & Continue source copy & 25 & $\{12\}$ & $D:[27,29,31,33,35,37,39,41,43,45]$\newline$x:48\quad \blank:61$ \\
48 & Descend below operator & 13 & $\{7\}$ & $D:49$ \\
49 & Write copied operator & 24 & $\{8\}$ & $\blank:50$ \\
50 & Copy last digit of $a$ & 25 & $\{9\}$ & $D:[51,52,53,54,55,56,57,58,59,60]$ \\
51 & Descend above $x$, digit 0 & 3 & $\{158\}$ & $x:28$ \\
52 & Descend above $x$, digit 1 & 4 & $\{145\}$ & $x:30$ \\
53 & Descend above $x$, digit 2 & 5 & $\{66\}$ & $x:32$ \\
54 & Descend above $x$, digit 3 & 6 & $\{10\}$ & $x:34$ \\
55 & Descend above $x$, digit 4 & 7 & $\{89\}$ & $x:36$ \\
56 & Descend above $x$, digit 5 & 8 & $\{130\}$ & $x:38$ \\
57 & Descend above $x$, digit 6 & 9 & $\{41\}$ & $x:40$ \\
58 & Descend above $x$, digit 7 & 10 & $\{147\}$ & $x:42$ \\
59 & Descend above $x$, digit 8 & 11 & $\{115\}$ & $x:44$ \\
60 & Descend above $x$, digit 9 & 12 & $\{53\}$ & $x:46$ \\
61 & Restart on copied row & 0 & $\{13\}$ & $D:1$ \\
62 & Bottom-row factor & 26 & $\{14\}$ & $D:[63,66,69,73,78,84,91,99,108,118]$ \\
63 & Erase $x$, bottom factor 0 & 27 & $\{111\}$ & $x:64$ \\
64 & First product digit, factor 0 & 37 & $\{62\}$ & $D:65$ \\
65 & Continue product by zero & 37 & $\{62\}$ & $D:65\quad \blank:129$ \\
66 & Erase $x$, bottom factor 1 & 28 & $\{78\}$ & $x:67$ \\
67 & First product digit, factor 1 & 38 & $\{24\}$ & $D:68$ \\
68 & Product by 1, carry 0 & 38 & $\{24\}$ & $D:68\quad \blank:129$ \\
69 & Erase $x$, bottom factor 2 & 29 & $\{69\}$ & $x:70$ \\
70 & First product digit, factor 2 & 39 & $\{70\}$ & $D:[71,71,71,71,71,72,72,72,72,72]$ \\
71 & Product by 2, carry 0 & 39 & $\{70\}$ & $D:[71,71,71,71,71,72,72,72,72,72]$\newline$\blank:129$ \\
72 & Product by 2, carry 1 & 40 & $\{143\}$ & $D:[71,71,71,71,71,72,72,72,72,72]$\newline$\blank:129$ \\
73 & Erase $x$, bottom factor 3 & 30 & $\{15\}$ & $x:74$ \\
74 & First product digit, factor 3 & 41 & $\{16\}$ & $D:[75,75,75,75,76,76,76,77,77,77]$ \\
75 & Product by 3, carry 0 & 41 & $\{16\}$ & $D:[75,75,75,75,76,76,76,77,77,77]$\newline$\blank:129$ \\
76 & Product by 3, carry 1 & 42 & $\{48\}$ & $D:[75,75,75,76,76,76,76,77,77,77]$\newline$\blank:129$ \\
77 & Product by 3, carry 2 & 43 & $\{155\}$ & $D:[75,75,75,76,76,76,77,77,77,77]$\newline$\blank:129$ \\
78 & Erase $x$, bottom factor 4 & 31 & $\{75\}$ & $x:79$ \\
79 & First product digit, factor 4 & 44 & $\{76\}$ & $D:[80,80,80,81,81,82,82,82,83,83]$ \\
80 & Product by 4, carry 0 & 44 & $\{76\}$ & $D:[80,80,80,81,81,82,82,82,83,83]$\newline$\blank:129$ \\
81 & Product by 4, carry 1 & 45 & $\{91\}$ & $D:[80,80,80,81,81,82,82,82,83,83]$\newline$\blank:129$ \\
82 & Product by 4, carry 2 & 46 & $\{77\}$ & $D:[80,80,81,81,81,82,82,83,83,83]$\newline$\blank:129$ \\
83 & Product by 4, carry 3 & 47 & $\{117\}$ & $D:[80,80,81,81,81,82,82,83,83,83]$\newline$\blank:129$ \\
84 & Erase $x$, bottom factor 5 & 32 & $\{43\}$ & $x:85$ \\
85 & First product digit, factor 5 & 48 & $\{44\}$ & $D:[86,86,87,87,88,88,89,89,90,90]$ \\
86 & Product by 5, carry 0 & 48 & $\{44\}$ & $D:[86,86,87,87,88,88,89,89,90,90]$\newline$\blank:129$ \\
87 & Product by 5, carry 1 & 49 & $\{153\}$ & $D:[86,86,87,87,88,88,89,89,90,90]$\newline$\blank:129$ \\
88 & Product by 5, carry 2 & 50 & $\{144\}$ & $D:[86,86,87,87,88,88,89,89,90,90]$\newline$\blank:129$ \\
89 & Product by 5, carry 3 & 51 & $\{45\}$ & $D:[86,86,87,87,88,88,89,89,90,90]$\newline$\blank:129$ \\
90 & Product by 5, carry 4 & 52 & $\{95\}$ & $D:[86,86,87,87,88,88,89,89,90,90]$\newline$\blank:129$ \\
91 & Erase $x$, bottom factor 6 & 33 & $\{107\}$ & $x:92$ \\
92 & First product digit, factor 6 & 53 & $\{108\}$ & $D:[93,93,94,94,95,96,96,97,97,98]$ \\
93 & Product by 6, carry 0 & 53 & $\{108\}$ & $D:[93,93,94,94,95,96,96,97,97,98]$\newline$\blank:129$ \\
94 & Product by 6, carry 1 & 54 & $\{171\}$ & $D:[93,93,94,94,95,96,96,97,97,98]$\newline$\blank:129$ \\
95 & Product by 6, carry 2 & 55 & $\{110\}$ & $D:[93,93,94,95,95,96,96,97,98,98]$\newline$\blank:129$ \\
96 & Product by 6, carry 3 & 56 & $\{174\}$ & $D:[93,93,94,95,95,96,96,97,98,98]$\newline$\blank:129$ \\
97 & Product by 6, carry 4 & 57 & $\{142\}$ & $D:[93,94,94,95,95,96,97,97,98,98]$\newline$\blank:129$ \\
98 & Product by 6, carry 5 & 58 & $\{109\}$ & $D:[93,94,94,95,95,96,97,97,98,98]$\newline$\blank:129$ \\
99 & Erase $x$, bottom factor 7 & 34 & $\{56\}$ & $x:100$ \\
100 & First product digit, factor 7 & 59 & $\{57\}$ & $D:[101,101,102,103,103,104,105,105,106,107]$ \\
101 & Product by 7, carry 0 & 59 & $\{57\}$ & $D:[101,101,102,103,103,104,105,105,106,107]$\newline$\blank:129$ \\
102 & Product by 7, carry 1 & 60 & $\{168\}$ & $D:[101,101,102,103,103,104,105,106,106,107]$\newline$\blank:129$ \\
103 & Product by 7, carry 2 & 61 & $\{127\}$ & $D:[101,101,102,103,104,104,105,106,106,107]$\newline$\blank:129$ \\
104 & Product by 7, carry 3 & 62 & $\{179\}$ & $D:[101,102,102,103,104,104,105,106,106,107]$\newline$\blank:129$ \\
105 & Product by 7, carry 4 & 63 & $\{59\}$ & $D:[101,102,102,103,104,104,105,106,107,107]$\newline$\blank:129$ \\
106 & Product by 7, carry 5 & 64 & $\{169\}$ & $D:[101,102,102,103,104,105,105,106,107,107]$\newline$\blank:129$ \\
107 & Product by 7, carry 6 & 65 & $\{58\}$ & $D:[101,102,103,103,104,105,105,106,107,107]$\newline$\blank:129$ \\
108 & Erase $x$, bottom factor 8 & 35 & $\{112\}$ & $x:109$ \\
109 & First product digit, factor 8 & 66 & $\{113\}$ & $D:[110,110,111,112,113,114,114,115,116,117]$ \\
110 & Product by 8, carry 0 & 66 & $\{113\}$ & $D:[110,110,111,112,113,114,114,115,116,117]$\newline$\blank:129$ \\
111 & Product by 8, carry 1 & 67 & $\{137\}$ & $D:[110,110,111,112,113,114,114,115,116,117]$\newline$\blank:129$ \\
112 & Product by 8, carry 2 & 68 & $\{121\}$ & $D:[110,111,111,112,113,114,115,115,116,117]$\newline$\blank:129$ \\
113 & Product by 8, carry 3 & 69 & $\{114\}$ & $D:[110,111,111,112,113,114,115,115,116,117]$\newline$\blank:129$ \\
114 & Product by 8, carry 4 & 70 & $\{132\}$ & $D:[110,111,112,112,113,114,115,116,116,117]$\newline$\blank:129$ \\
115 & Product by 8, carry 5 & 71 & $\{149\}$ & $D:[110,111,112,112,113,114,115,116,116,117]$\newline$\blank:129$ \\
116 & Product by 8, carry 6 & 72 & $\{120\}$ & $D:[110,111,112,113,113,114,115,116,117,117]$\newline$\blank:129$ \\
117 & Product by 8, carry 7 & 73 & $\{136\}$ & $D:[110,111,112,113,113,114,115,116,117,117]$\newline$\blank:129$ \\
118 & Erase $x$, bottom factor 9 & 36 & $\{81\}$ & $x:119$ \\
119 & First product digit, factor 9 & 74 & $\{82\}$ & $D:[120,120,121,122,123,124,125,126,127,128]$ \\
120 & Product by 9, carry 0 & 74 & $\{82\}$ & $D:[120,120,121,122,123,124,125,126,127,128]$\newline$\blank:129$ \\
121 & Product by 9, carry 1 & 75 & $\{151\}$ & $D:[120,121,121,122,123,124,125,126,127,128]$\newline$\blank:129$ \\
122 & Product by 9, carry 2 & 76 & $\{177\}$ & $D:[120,121,122,122,123,124,125,126,127,128]$\newline$\blank:129$ \\
123 & Product by 9, carry 3 & 77 & $\{154\}$ & $D:[120,121,122,123,123,124,125,126,127,128]$\newline$\blank:129$ \\
124 & Product by 9, carry 4 & 78 & $\{176\}$ & $D:[120,121,122,123,124,124,125,126,127,128]$\newline$\blank:129$ \\
125 & Product by 9, carry 5 & 79 & $\{83\}$ & $D:[120,121,122,123,124,125,125,126,127,128]$\newline$\blank:129$ \\
126 & Product by 9, carry 6 & 80 & $\{156\}$ & $D:[120,121,122,123,124,125,126,126,127,128]$\newline$\blank:129$ \\
127 & Product by 9, carry 7 & 81 & $\{152\}$ & $D:[120,121,122,123,124,125,126,127,127,128]$\newline$\blank:129$ \\
128 & Product by 9, carry 8 & 82 & $\{98\}$ & $D:[120,121,122,123,124,125,126,127,128,128]$\newline$\blank:129$ \\
129 & Scan completed product & 83 & $\{17\}$ & $D:129\quad \blank:130$ \\
130 & Inspect row above & 84 & $\{18\}$ & $D:131\quad \blank:155$ \\
131 & Cross upper operator & 85 & $\{19\}$ & $x:132$ \\
132 & First upper factor digit & 85 & $\{20\}$ & $D:133$ \\
133 & Scan upper factor prefix & 85 & $\{20\}$ & $D:133\quad \blank:134$ \\
134 & Select upper factor & 26 & $\{21\}$ & $D:[135,137,139,141,143,145,147,149,151,153]$ \\
135 & First prefix erasure, factor 0 & 27 & $\{60\}$ & $D:136$ \\
136 & Later prefix erasure, factor 0 & 27 & $\{61\}$ & $D:136\quad x:64$ \\
137 & First prefix erasure, factor 1 & 28 & $\{22\}$ & $D:138$ \\
138 & Later prefix erasure, factor 1 & 28 & $\{23\}$ & $D:138\quad x:67$ \\
139 & First prefix erasure, factor 2 & 29 & $\{71\}$ & $D:140$ \\
140 & Later prefix erasure, factor 2 & 29 & $\{72\}$ & $D:140\quad x:70$ \\
141 & First prefix erasure, factor 3 & 30 & $\{46\}$ & $D:142$ \\
142 & Later prefix erasure, factor 3 & 30 & $\{47\}$ & $D:142\quad x:74$ \\
143 & First prefix erasure, factor 4 & 31 & $\{92\}$ & $D:144$ \\
144 & Later prefix erasure, factor 4 & 31 & $\{93\}$ & $D:144\quad x:79$ \\
145 & First prefix erasure, factor 5 & 32 & $\{84\}$ & $D:146$ \\
146 & Later prefix erasure, factor 5 & 32 & $\{85\}$ & $D:146\quad x:85$ \\
147 & First prefix erasure, factor 6 & 33 & $\{140\}$ & $D:148$ \\
148 & Later prefix erasure, factor 6 & 33 & $\{141\}$ & $D:148\quad x:92$ \\
149 & First prefix erasure, factor 7 & 34 & $\{160\}$ & $D:150$ \\
150 & Later prefix erasure, factor 7 & 34 & $\{161\}$ & $D:150\quad x:100$ \\
151 & First prefix erasure, factor 8 & 35 & $\{118\}$ & $D:152$ \\
152 & Later prefix erasure, factor 8 & 35 & $\{119\}$ & $D:152\quad x:109$ \\
153 & First prefix erasure, factor 9 & 36 & $\{96\}$ & $D:154$ \\
154 & Later prefix erasure, factor 9 & 36 & $\{97\}$ & $D:154\quad x:119$ \\
155 & Start summation & 86 & $\{25\}$ & $\blank:156$ \\
156 & First initial left digit & 87 & $\{26\}$ & $D:157$ \\
157 & Continue initial left scan & 87 & $\{26\}$ & $D:157\quad \blank:158$ \\
158 & Inspect fresh lower row & 88 & $\{27\}$ & $0{:}8:159\quad \blank:205$ \\
159 & First digit after leading cell & 89 & $\{28\}$ & $D:160$ \\
160 & Scan lower product & 89 & $\{28\}$ & $D:160\quad \blank:161$ \\
161 & First upper alignment digit & 104 & $\{29\}$ & $D:162$ \\
162 & Continue upper alignment & 104 & $\{30\}$ & $D:162\quad \blank:163$ \\
163 & Read upper units digit & 90 & $\{31\}$ & $D:[164,165,166,167,168,169,170,171,172,173]$ \\
164 & First lower update, value 0 & 92 & $\{63\}$ & $\blank:174$ \\
165 & First lower update, value 1 & 93 & $\{99\}$ & $\blank:174$ \\
166 & First lower update, value 2 & 94 & $\{138\}$ & $\blank:174$ \\
167 & First lower update, value 3 & 95 & $\{32\}$ & $\blank:174$ \\
168 & First lower update, value 4 & 96 & $\{73\}$ & $\blank:174$ \\
169 & First lower update, value 5 & 97 & $\{133\}$ & $\blank:174$ \\
170 & First lower update, value 6 & 98 & $\{94\}$ & $\blank:174$ \\
171 & First lower update, value 7 & 99 & $\{162\}$ & $\blank:174$ \\
172 & First lower update, value 8 & 100 & $\{49\}$ & $\blank:174$ \\
173 & First lower update, value 9 & 101 & $\{146\}$ & $\blank:174$ \\
174 & Return after lower blank & 91 & $\{33\}$ & $\blank:175$ \\
175 & Read next upper digit & 90 & $\{34\}$ & $D:[178,179,180,181,182,183,184,185,186,187]$ \\
176 & Return with carry zero & 91 & $\{36\}$ & $\blank:177$ \\
177 & Read upper, carry zero & 90 & $\{37\}$ & $D:[178,179,180,181,182,183,184,185,186,187]$\newline$\blank:200$ \\
178 & Add 0 below & 92 & $\{35\}$ & $D:176\quad \blank:174$ \\
179 & Add 1 below & 93 & $\{50\}$ & $D:[176,176,176,176,176,176,176,176,176,188]$\newline$\blank:174$ \\
180 & Add 2 below & 94 & $\{122\}$ & $D:[176,176,176,176,176,176,176,176,188,188]$\newline$\blank:174$ \\
181 & Add 3 below & 95 & $\{86\}$ & $D:[176,176,176,176,176,176,176,188,188,188]$\newline$\blank:174$ \\
182 & Add 4 below & 96 & $\{74\}$ & $D:[176,176,176,176,176,176,188,188,188,188]$\newline$\blank:174$ \\
183 & Add 5 below & 97 & $\{139\}$ & $D:[176,176,176,176,176,188,188,188,188,188]$\newline$\blank:174$ \\
184 & Add 6 below & 98 & $\{157\}$ & $D:[176,176,176,176,188,188,188,188,188,188]$\newline$\blank:174$ \\
185 & Add 7 below & 99 & $\{175\}$ & $D:[176,176,176,188,188,188,188,188,188,188]$\newline$\blank:174$ \\
186 & Add 8 below & 100 & $\{100\}$ & $D:[176,176,188,188,188,188,188,188,188,188]$\newline$\blank:174$ \\
187 & Add 9 below & 101 & $\{126\}$ & $D:[176,188,188,188,188,188,188,188,188,188]$\newline$\blank:174$ \\
188 & Return with carry one & 103 & $\{101\}$ & $\blank:189$ \\
189 & Read upper, carry one & 102 & $\{102\}$ & $D:[190,191,192,193,194,195,196,197,198,199]$\newline$\blank:202$ \\
190 & Add 1 inside lower block & 93 & $\{50\}$ & $D:[176,176,176,176,176,176,176,176,176,188]$ \\
191 & Add 2 inside lower block & 94 & $\{122\}$ & $D:[176,176,176,176,176,176,176,176,188,188]$ \\
192 & Add 3 inside lower block & 95 & $\{86\}$ & $D:[176,176,176,176,176,176,176,188,188,188]$ \\
193 & Add 4 inside lower block & 96 & $\{74\}$ & $D:[176,176,176,176,176,176,188,188,188,188]$ \\
194 & Add 5 inside lower block & 97 & $\{139\}$ & $D:[176,176,176,176,176,188,188,188,188,188]$ \\
195 & Add 6 inside lower block & 98 & $\{157\}$ & $D:[176,176,176,176,188,188,188,188,188,188]$ \\
196 & Add 7 inside lower block & 99 & $\{175\}$ & $D:[176,176,176,188,188,188,188,188,188,188]$ \\
197 & Add 8 inside lower block & 100 & $\{100\}$ & $D:[176,176,188,188,188,188,188,188,188,188]$ \\
198 & Add 9 inside lower block & 101 & $\{126\}$ & $D:[176,188,188,188,188,188,188,188,188,188]$ \\
199 & Add 10 inside lower block & 105 & $\{180\}$ & $D:188$ \\
200 & Preserve leading digit & 87 & $\{26\}$ & $0{:}8:201$ \\
201 & Blank left of merged row & 87 & $\{26,106\}$ & $\blank:158$ \\
202 & Increment leading digit & 93 & $\{103\}$ & $0{:}8:203$ \\
203 & Return after final carry & 91 & $\{104\}$ & $\blank:204$ \\
204 & Step down left of sum & 90 & $\{105\}$ & $\blank:201$ \\
205 & Emit final stop & 106 & $\{38\}$ & $\blank:206$ \\
206 & Halted & 107 & $\{181\}$ & $\text{--}$ \\
\end{longtable}
\endgroup

\bigskip
\sectiontitle{5 \enspace EQUIVALENCE PROOF}{Equivalence Proof}\label{sec:common5}
\begin{proof}[Proof of Theorem~\ref{thm:main}]
Initially both executions have the same grid and head; phase $0$ relates
the two initial control states. Suppose the relation holds before a step.
Lemma~\ref{lem:coverage} supplies a permitted read symbol $x$.
Equal grids and heads give the same read symbol to both controllers.
Equation~(4) gives the same write and motion and relates the successor
states through phase $j$. The next grids and heads therefore agree.

Induction establishes equality of every finite execution prefix.
The halt condition permits termination exactly on entry to phase $206$.
The expert terminates by the finite copy, product, and merge arguments;
the extracted controller therefore terminates on the same step.
The complete PTOs and final configurations agree.
Lemma~\ref{lem:sum} supplies the final padded product and head position.
\end{proof}

\textbf{Undefined Cells.}
The partial controller leaves $182\cdot12-1{,}271=913$ entries undefined.
None occurs on a valid multiplication: protocol coverage and (4) supply
a defined entry at every possible execution step.
For example, controller state $0$ scans the left operand and never reads
a blank, whereas state $1$ scans the right operand and handles the right
blank. Both states correspond to expert state $0$; merging these semantic
roles during inspection would incorrectly demand a blank entry at controller
state $0$. The phase relation retains the necessary context.

\bigskip
\sectiontitle{6 \enspace VERIFICATION AND ARTIFACTS}{Verification and Artifacts}\label{sec:common6}
\textbf{Verification.}
Run \code{python3 verify\_equivalence.py {-}{-}quick} from this directory.
The checker verifies fingerprints, CSV/JSON equality, equality of the expert
export with \code{mult\_grid(10)}, and every obligation in (4).
The checked phase relation matches every row of the certificate.
Omitting \code{{-}{-}quick} additionally replays all $2{,}000$ extraction inputs
and the $100$ saved pairs of $100$-digit operands on separate autonomous grids.
The $180{,}510$ extraction actions and $15{,}114{,}148$ long-input actions
agree exactly; final grids, heads, and padded arithmetic products also agree.
The invariants and induction establish the unbounded claim.

\textbf{Artifact Identity.}
The title page gives the repository-relative case directory. This proof directory
contains the controller and expert exports, reference code, protocol, checker,
and saved verification report. The verification script stores or reads SHA-256 checksums for these artifacts. A checksum identifies the exact bytes
of a file; the checker compares the stored and computed values to detect changes.
These checks bind the certificate to the supplied exports. The invariants,
protocol coverage, and local transition identity establish equivalence.
\end{document}
